% Part: incompleteness
% Chapter: arithmetization-syntax
% Section: coding-terms

\documentclass[../../../include/open-logic-section]{subfiles}

\begin{document}

\olfileid{inc}{art}{trm}
\olsection{పదాల సంకేతీకరణ}

\begin{explain}
పదం అంటే ఒక నిర్దిష్ట రకమైన సంకేతాల క్రమమే: పదాల నిర్మాణ నియమాల ప్రకారం
స్థిరసంకేతాలు, చరాల నుంచి
దాన్ని ఆగమనాత్మకంగా నిర్మిస్తాం. సంకేతాల కోసం ఒక సంకేతీకరణ పద్ధతితో పాటు
సంఖ్యల క్రమాలను సంకేతీకరించే పద్ధతిని ఉపయోగించి సంకేతాల క్రమాలను సంఖ్యలుగా
సంకేతీకరించవచ్చు కనుక, పదాలకు గ్యోడెల్ సంఖ్యలను కేటాయించడం కష్టం కాదు.
సరిగ్గా నిర్మించిన పదానికి గ్యోడెల్ సంఖ్యగా ఉండే సంఖ్యకు గల ధర్మం గణనీయమని,
నిజానికి ఆదిమ పునరావృత్తమని చూపడమే అసలు సవాలు.
\end{explain}

\tetoken{చరాలు}{!!^{variable}s}, \tetoken{స్థిరసంకేతాలు}{!!{constant}s}
అత్యంత సరళమైన పదాలు; $x$ అటువంటి పదానికి గ్యోడెల్ సంఖ్య కాదా అని పరీక్షించడం
సులభం: ఏదో~$i$కి $x$ అనేది $\Gn{\Obj v_i}$ అయితే $\fn{Var}(x)$ నిలుస్తుంది.
మరో మాటలో, $x$ పొడవు~$1$ గల క్రమం, దాని ఏకైక మూలకం $(x)_0$ ఏదో
\tetoken{చరం}{!!{variable}}~$\Obj v_i$ సంకేతసంఖ్య; అంటే ఏదో~$i$కి
$x$ అనేది $\tuple{\tuple{1,i}}$. ఇదే విధంగా, ఏదో~$i$కి $x$ అనేది $\Gn{\Obj c_i}$
అయితే $\fn{Const}(x)$ నిలుస్తుంది. అటువంటి $i$ ఉంటే అది $<x$ కావాలి కనుక,
ఈ రెండు సంబంధాలూ ఆదిమ పునరావృత్తం:
\begin{align*}
  \fn{Var}(x) & \defiff \bexists{i<x}{x = \tuple{\tuple{1, i}}}\\
  \fn{Const}(x) & \defiff \bexists{i<x}{x = \tuple{\tuple{2, i}}}
\end{align*}

\begin{prop}
\ollabel{prop:term-primrec}
$x$ వరుసగా ఒక పదానికి లేదా ఒక మూసిన పదానికి గ్యోడెల్ సంఖ్య అయితేనే నిలిచే
$\fn{Term}(x)$, $\fn{ClTerm}(x)$ సంబంధాలు ఆదిమ పునరావృత్తం.
\end{prop}

\begin{proof}
సంకేతాల క్రమం~$s$ ఒక పదం కావడానికి అవసరమైనదీ, సరిపడేదీ ఏమిటంటే, పదాల
నిర్మాణ నియమాల ప్రకారం \tetoken{స్థిరసంకేతాలు}{!!{constant}s},
\tetoken{చరాల}{!!{variable}s} నుంచి $s$ ఎలా ఏర్పడిందో నమోదు చేసే
$s_0$, \dots, $s_{k-1}=s$ అనే పదాల క్రమం ఉండటం. అటువంటి ఊహాత్మక నిర్మాణ
క్రమం నిర్మాణ నియమాలను పాటిస్తుందని చెప్పాలంటే ప్రతి $i<k$కి క్రింది వాటిలో
ఒకటి జరగాలి:
\begin{enumerate}
\item $s_i$ ఒక \tetoken{చరం}{!!a{variable}}~$\Obj v_j$, లేదా
\item $s_i$ ఒక \tetoken{స్థిరసంకేతం}{!!a{constant}}~$\Obj c_j$, లేదా
\item $s_i$ అనేది $i$-వ స్థానానికి ముందు కనిపించే $n$ పదాలు
  $t_1$, \dots, $t_n$ నుంచి, $n$-స్థాన
  \tetoken{ప్రమేయ సంకేతం}{!!{function}}~$\Obj f^n_j$ ఉపయోగించి నిర్మితమైంది.
\end{enumerate}
గ్యోడెల్ సంఖ్యలపై దీనికి అనురూపమైన సంబంధం ఆదిమ పునరావృత్తమని చూపడానికి,
ఈ షరతును ఆదిమ పునరావృత్త ప్రమేయాలు, సంబంధాలు, పరిమిత పరిమాణీకరణతో
ఆదిమ పునరావృత్తంగా వ్యక్తీకరించాలి.

$y$ అనేది $s_0$, \dots, $s_{k-1}$ క్రమాన్ని సంకేతీకరించే సంఖ్య అని,
అంటే $y=\tuple{\Gn{s_0},\dots,\Gn{s_{k-1}}}$ అని అనుకుందాం. ప్రతి $i<k$కి
క్రింది వాటిలో ఒకటి జరిగితేనే, ఇంకా $(y)_{k-1}=x$ అయితేనే, అది గ్యోడెల్
సంఖ్య~$x$ గల పదానికి నిర్మాణ క్రమాన్ని సంకేతీకరిస్తుంది:
\begin{enumerate}
\item $\fn{Var}((y)_i)$, లేదా
\item $\fn{Const}((y)_i)$, లేదా
\item $n$, $j$ అనే సంఖ్యలూ, $z=\tuple{z_1,\dots,z_n}$ అనే సంఖ్యా ఉన్నాయి;
  ప్రతి $z_l$ ఏదో $i'<i$కి $(y)_{i'}$తో సమానం, అలాగే
\[
(y)_i = \Gn{\Obj f^n_j(} \concat \fn{flatten}(z) \concat \Gn{)},
\]
\end{enumerate}
($\fn{flatten}(z)$ ప్రమేయం $\tuple{\Gn{t_1},\dots,\Gn{t_n}}$ క్రమాన్ని
$\Gn{t_1,\dots,t_n}$గా మారుస్తుంది; అది ఆదిమ పునరావృత్తం.)
\sourcecorrection{OLTEART-002}{ఆంగ్ల మూలంలోని నిర్మాణ షరతు (3)లో $f^n_j$ను ఉపయోగించినప్పటికీ, ఉనికిపరిమాణీకరణలో $j$ను పేర్కొనలేదు. వెంటనే వచ్చే వాక్యం $j$ కూడా $y$కంటే తక్కువని స్పష్టంగా చెబుతుంది; అందుకు అనుగుణంగా ఇక్కడ $j$ను ఉనికిపరిమాణీకరించాం.}

(3)లోని సూచికలు $j$, $n$, పదాలు $t_l$ యొక్క గ్యోడెల్ సంఖ్యలు $z_l$,
క్రమం $\tuple{z_1,\dots,z_n}$ యొక్క సంకేతసంఖ్య~$z$ అన్నీ $y$కంటే తక్కువ.
పైన $k$ స్థానంలో $\len{y}$ను పెట్టవచ్చు. కాబట్టి ``$y$ అనేది గ్యోడెల్
సంఖ్య~$x$ గల పదానికి నిర్మాణ క్రమపు సంకేతసంఖ్య'' అనే వాక్యాన్ని, ఈ సంబంధం
ఆదిమ పునరావృత్తమని చూపే విధంగా వ్యక్తీకరించవచ్చు.

ఇప్పుడు $y$పై ఒక ఆదిమ పునరావృత్త పరిమితి ఉందని నిర్ధారించుకుంటే చాలు.
$x$ ఒక పదానికి గ్యోడెల్ సంఖ్య అయితే, దానికి గరిష్ఠంగా $\len{x}$ పదాలు గల
నిర్మాణ క్రమం ఉండాలి (ఎందుకంటే $s$ యొక్క నిర్మాణ క్రమంలోని ప్రతి పదం
$s$లో ఏదో ఒక స్థానంలో మొదలవ్వాలి; రెండు ఉపపదాలు ఒకే స్థానంలో మొదలవ్వలేవు).
$s$ యొక్క ప్రతి ఉపపదపు గ్యోడెల్ సంఖ్య సహజంగానే $\le x$. అందువల్ల,
$k=\len{x}$ అయినప్పుడు సంకేతసంఖ్య $\le p_{k-1}^{k(x+1)}$ గల నిర్మాణ
క్రమం ఎప్పుడూ ఉంటుంది.

$\fn{ClTerm}$ కోసం, \tetoken{చరాల}{!!{variable}s} విభాగాన్ని వదిలివేస్తే చాలు.
\end{proof}

\begin{prob}
$\tuple{\Gn{t_1},\dots,\Gn{t_n}}$ క్రమాన్ని $\Gn{t_1,\dots,t_n}$గా మార్చే
$\fn{flatten}(z)$ ప్రమేయం ఆదిమ పునరావృత్తమని చూపండి.
\end{prob}

\begin{prop}\ollabel{prop:num-primrec}
  $\fn{num}(n)=\Gn{\num{n}}$ ప్రమేయం ఆదిమ పునరావృత్తం.
\end{prop}

\begin{proof}
  $\fn{num}(n)$ను ఆదిమ పునరావృత్తం ద్వారా ఇలా నిర్వచిస్తాం:
  \begin{align*}
    \fn{num}(0) & = \Gn{\Obj 0}\\
    \fn{num}(n+1) & = \Gn{\prime(} \concat \fn{num}(n) \concat \Gn{)}.
  \end{align*}
\end{proof}

\end{document}
